% Les quatre étapes d'une synthèse organique (L'essentiel, chapitre 12)
\begin{tikzpicture}[x=1cm,y=1cm,
  etape/.style={draw=otNavy,fill=otNavyTint,line width=0.6pt,rounded corners=1.4mm,minimum width=26mm,minimum height=11mm,align=center,font=\otMedium\fontsize{8.5}{10}\selectfont,text=otNavy}]
\node[etape] (a) at (0,0) {transformation};
\node[etape] (b) at (3.4,0) {isolement};
\node[etape] (c) at (6.8,0) {purification};
\node[etape] (d) at (10.2,0) {identification};
\foreach \u/\v in {a/b,b/c,c/d}{\draw[pcFleche] (\u.east) -- (\v.west);}
\node[pcLegende,anchor=north,align=center] at (0,-0.75) {chauffage à reflux};
\node[pcLegende,anchor=north,align=center] at (3.4,-0.75) {filtration, décantation,\\extraction};
\node[pcLegende,anchor=north,align=center] at (6.8,-0.75) {recristallisation,\\distillation, lavages};
\node[pcLegende,anchor=north,align=center] at (10.2,-0.75) {CCM, température\\de fusion, infrarouge};
\end{tikzpicture}
\par\smallskip\legende{Chaque étape a son objectif : transformer, récupérer le produit brut, le purifier, puis vérifier son identité et sa pureté.}
